In Cosmology, Cepheid variables are called standard candles as they can be used to estimate distances. Edwin Hubble used these standard candles to settle the Shapley--Curtis debate by measuring the distance to the Andromeda Galaxy. Cepheid variables are variable stars whose brightness vary periodically. Henrietta Leavitt discovered the relation between the periodicity and average luminosity of these stars, and this relation can be used to derive the equation linking Absolute magnitude and Period. For Cepheid variables, luminosity $L$ varies with period $P$ as $L\propto P^{2.5}$, where $P$ is in days.
\begin{enumerate}
\item Find the relation between Absolute Magnitude $M$ and $P$ for this type of variable stars.
\item You might notice that it is a linear equation between $M$ and a function of $P$. If the unknown constant of the given equation is $-1.5\,\mathrm{mag}$, then find the distance to the given Cepheid Variable.
\end{enumerate}
\ioaafig{2024_TH_T04-f1.pdf}
